\documentclass[%
sor,
 jor,
 amsmath,amssymb,
 reprint,
]{revtex4-2}
\usepackage{float}
\usepackage{caption}
\usepackage{amsmath, amssymb, amsfonts}
\usepackage{siunitx}
\usepackage{graphicx}
\usepackage{circuitikz}



\begin{document}
\title{PH3244\\Experiment - \\ }
\author{Aumshree P. Shah\\20231059}
\altaffiliation{\color{red}aumshree.pinkalbenshah@students.iiserpune.ac.in}
\date{\today}
\begin{abstract}
	ABSTRACT,  write at last, add comments her
\\WRITE EXP NO AND NAME
\\theory
\\check a llook at prcedure
\\ TILL HERE MAKE SURE YOU CLEAR THE TOP TWO AND THEM MOVE FORWARDS
\\ TILL HERE MAKE SURE TO FINALIZE THE RESULT FROM CALC
\\error and prop formula with ref
\\do the calc and result
\centering
\end{abstract}
\maketitle
\tableofcontents
\vspace*{\fill}
\pagebreak
\pagebreak
\section{Theory and Procedure}
	\subsection{Theory\footnote{Please refer to [\cite{ph3244repo}] for a full examination of theory, only the important results will be discussed here.}}

\subsection{Procedure}Refer to the manual [\cite{ph3244repo}] for a look at procedure, although the required pages will be attached at the back for sumbition along with the signed report.

\section{Observation}
\subsection{For constant $V_{\text{DS}}$}
\begin{table}[H]
\begin{minipage}{0.48\linewidth}
	\begin{center}
\begin{tabular}{|c|c|}
\hline
$V_{\text{GS}}$  (V) & $I_D$ (mA) \\
\hline
0.000 & 7.17 \\
\hline
\end{tabular}
\caption{$V_{\text{DS}} = -2.65\si{\volt}$}
\end{center}
\end{minipage}
\begin{minipage}{0.48\linewidth}
	\begin{center}


\begin{tabular}{|c|c|}
\hline
$V_{\text{GS}}$  (V) & $I_D$ (mA) \\
\hline
0.000 & 8.58 \\
\hline
\end{tabular}
\caption{$V_{\text{DS}} = -5\si{\volt}$}
\end{center}
\end{minipage}
\hfill
\end{table}
\begin{center}
	Data taken on 6 Feb 2026 
\end{center}

\section{Calculations, Analysis and Characteristics}
\subsection{Sources of Error}
\subsubsection{Systematic Errors}
\begin{enumerate}
	\item A
\end{enumerate}
\subsubsection{Random Errors}
\begin{enumerate}
	\item Uncertanity in instruments
\end{enumerate}
\subsection{Error Propogation}
To calculate the standard deviation from values we have measured with some least count, and the values we have applied digitally, we use the formula [\cite{tewari_uncertainty}]: 

\begin{equation}
\sigma = \frac{\text{Least Count}}{\sqrt{12}}
\end{equation}
and to calculate the standard deviation for values we have applied on analog instrument with some least count, we use the formula [\cite{tewari_uncertainty}]: 
\begin{equation}
\sigma = \frac{\text{Least Count}}{\sqrt{24}}
\end{equation}
ERRROR PROPEOGASTION AND CITE:
[\cite{taylor_error_analysis}] that: 
\begin{equation}
	\Delta \mu = \mu \sqrt {\left( \frac {\Delta g_m }{g_m} \right)^2 + \left( \frac {\Delta r_d}{r_d} \right)^2}
\end{equation}

\subsection{Plots and Calculation\footnote{Refer to [\cite{ph3244repo}] for the code used to plot and determine the slope.}}
\begin{figure}[H]
\begin{minipage}{0.48\linewidth}
\begin{center}
	%\includegraphics[width=1\linewidth]{obs1_nofit}
	\caption{$I_d$ vs $V_{GS}$ from Table - I, II}
\end{center}

\end{minipage}
\begin{minipage}{0.48\linewidth}
\begin{center}
	%\includegraphics[width=1\linewidth]{obs2_nofit}
	\caption{$I_d$ vs $V_{DS}$ from Table - III, IV, V}
\end{center}
\end{minipage}
\end{figure}






\section{Result and Conclusion}

\cite{ph3244repo}
\cite{tewari_uncertainty}
\cite{taylor_error_analysis}
\cite{standard_error_wiki}

\vspace{1cm}
\hrule
\vspace{1cm}
\appendix
\section{Appendis}
theo val.\\
wat cam be o,[rpved

\bibliographystyle{apsrev4-2}
\bibliography{refrences}

\end{document}
